\label{sec:related_work}
\paragraph{Zero-shot training of retrievers.}
Recent neural retrievers~\cite{karpukhin-etal-2020-dense,lee-etal-2019-latent,colbert} show their superiority over term-based retrievers (e.g., BM25; \citealt{robertson2009probabilistic}) across domains when training data is abundant~\cite{luo2022improving,asai2021one,petroni-etal-2021-kilt}.
{Due to the high cost of annotating retrieval datasets for new target tasks, improving neural retrievers in zero-shot settings is an active area of study. 
The first line of work uses purely unsupervised approaches, such as pre-training neural retrievers on {unlabeled} data (e.g., Contriever; \citealt{izacard2022unsupervised}).
}
The second line of work trains a single retrieval system on large-scale supervised datasets such as MS MARCO~\cite{msmarco} and then performs a transfer to new datasets~\cite{izacard2022unsupervised,colbert,nogueira-etal-2020-document,chen2021salient}, which often struggle with tasks unlike those used for training~\cite{promptagator}. 
To address this, the third paradigm trains customized retrievers for each task using unlabeled corpora, leveraging another model to automatically generate training data~\cite{wang-etal-2022-gpl}. 
Concurrent to our work, \citet{promptagator} use task-specific templates and few-shot samples to automatically generate in-domain training queries given randomly sampled documents from the target corpus using FLAN~\cite{wei2022finetuned}. 
It often entails running massive LMs and training separate retrievers, resulting in slow and costly adaptation. 
\begin{table*}[t!]
\renewcommand{\arraystretch}{1.2}
\setlength{\tabcolsep}{2pt}
\footnotesize
    \centering
    \begin{tabular}{ll}%|ccccccc}
\toprule
\textbf{Dataset} & \textbf{Instruction} \\\midrule
% {\bf Training} & \\
NQ & \texttt{Retrieve a \hlc[cyan!30]{Wikipedia} \hlc[green!30]{paragraph} that \hlc[pink]{answers this question}.}  \\
Med Simple & \texttt{Your task is to find a \hlc[pink]{simplified} \hlc[green!30]{paragraph} of this paragraph from a \hlc[cyan!30]{medical paper.} }\\
% \hline
% {\bf Evaluation} & \\
QReCC & \texttt{Find a \hlc[green!30]{dialogue response} from  \hlc[cyan!30]{dialogue history} to \hlc[pink]{answer the user's question}. }  \\
Arguana & \texttt{Retrieve a \hlc[green!30]{paragraph} from \hlc[cyan!30]{an argument website} that \hlc[pink]{argues against the following argument}.}  \\
SciFact & \texttt{Find a \hlc[green!30]{sentence} from a \hlc[cyan!30]{scientific paper} to check if \hlc[pink]{the statement is correct or not}. } \\
MultiLexSum & \texttt{I want to find the \hlc[green!30]{one-sentence} \hlc[pink]{summary of this} \hlc[cyan!30]{legal case}.} \\
\bottomrule
 \end{tabular}
    \caption{Example instructions for natural questions (NQ; \citealt{47761}), medical text simplification (Med Simple;~\citealt{devaraj-etal-2021-paragraph}), QReCC~\cite{qrecc}, Arguana~\cite{wachsmuth-etal-2018-retrieval}, SciFact~\cite{wadden-etal-2020-fact} and MultiLexSum~\cite{shen2022multilexsum}. Each instruction defines \hlc[pink]{\emph{intent}}, \hlc[cyan!30]{\emph{domain}} and \hlc[green!30]{\emph{unit}}. 
    }\label{tab:instructions}
\end{table*}
\paragraph{Instruction tuning.}
Training large language models (LLMs) with instructions or demonstrations on many tasks has proven very effective for zero- or few-shot transfer in a variety of settings~\cite{wei2022finetuned,sanh2022multitask,instruct_gpt,min-etal-2022-metaicl,wang2022benchmarking,mishra-etal-2022-cross,flan_palm}. 
However, such instruction tuning has not been systematically explored in retrieval for several reasons.
First, large-scale instruction-annotated datasets~\cite{bach2022promptsource,wang2022benchmarking} do not include retrieval tasks.
Second, successful instruction-following LLMs are encoder-decoder or decoder-only models with tens of billions of parameters, which are difficult to be adapted for retrieval tasks, as we typically need to encode millions of documents. 

\paragraph{Retrieval with descriptions.}
{
To incorporate more fine-grained information needs, retrieval with descriptions or narratives (e.g., TREC 2004 Robust Track; \citealt{96071}) has been studied. 
Descriptions or narratives are more detailed natural language explanations that describe the information needs (i.e., desirable documents) for each query mainly for query disambiguation. 
However, early work shows that concatenating descriptions or narratives only marginally outperforms the baselines with titles only~\cite{walker1998okapi,10.1145/3404835.3462812}.
More recent work~\cite{DBLP:conf/sigir/DaiC19,10.1145/3366423.3380258} suggests that powerful encoders such as BERT~\cite{devlin-etal-2019-bert} could better incorporate rich linguistic context and boost performance.  }

